\section*{Chapter \ref{ch:g12:polymers} --- Polymers}

\begin{solution}{exo:g12:polymers:1}
Propene, \ce{CH2=CH-CH3}.
\end{solution}

\begin{solution}{exo:g12:polymers:2}
\ce{-CF2-CF2-}; an \omterm{def:g12:polymers:addition-polymer}{addition polymer}: the \omterm{def:g12:polymers:polymer}{monomer} has a \ce{C=C} bond and
no other product forms.
\end{solution}

\begin{solution}{exo:g12:polymers:3}
Addition: polyethene, PVC, polystyrene. Condensation: PET (a polyester),
nylon (a polyamide).
\end{solution}

\begin{solution}{exo:g12:polymers:4}
Cellulose, starch, rubber from latex, proteins. Nylon and PVC are
synthetic.
\end{solution}

\begin{solution}{exo:g12:polymers:5}
Each link uses one group of each \omterm{def:g12:polymers:polymer}{monomer}; with two groups, a \omterm{def:g12:polymers:polymer}{monomer}
can bond on both sides and the chain keeps growing. A \omterm{def:g12:polymers:polymer}{monomer} with one
group would stop the chain.
\end{solution}

\begin{solution}{exo:g12:polymers:6}
$M(\ce{C2H3Cl}) = 24.0 + 3.0 + 35.5 = \qty{62.5}{g/mol}$;
$n = 125000/62.5 = 2000$.
\end{solution}

\begin{solution}{exo:g12:polymers:7}
$M(\ce{C8H8}) = \qty{104.0}{g/mol}$; $M = 2500 \times 104.0 =
\qty{2.60e5}{g/mol}$.
\end{solution}

\begin{solution}{exo:g12:polymers:8}
The acid group of one \omterm{def:g7:atoms-and-molecules:molecule}{molecule} forms an \omterm{def:g11:functional-groups:carboxylic-acid}{ester} with the \omterm{def:g11:functional-groups:alcohol}{alcohol} group of
the next, which still has a free acid group: the chain grows. \omterm{def:g12:polymers:repeat-unit}{Repeat
unit} \ce{-O-CH(CH3)-CO-}; water is released at each link.
\end{solution}

\begin{solution}{exo:g12:polymers:9}
Bag films are made of branched chains, which cannot pack: the material
is amorphous, soft and transparent. Milk bottles are made of nearly
linear chains, which pack into crystalline regions: denser, stiffer, and
the crystalline regions scatter light.
\end{solution}

\begin{solution}{exo:g12:polymers:10}
Their chains are joined by many covalent cross-links. Heating cannot
make the chains slide; it breaks bonds and destroys the material instead
of melting it.
\end{solution}

\begin{solution}{exo:g12:polymers:11}
$M(\ce{C12H22N2O2}) = 144.0 + 22.0 + 28.0 + 32.0 = \qty{226.0}{g/mol}$;
$n = 22600/226.0 = 100$.
\end{solution}

\begin{solution}{exo:g12:polymers:12}
\[
  \ce{H2N(CH2)6NH2 + HOOC(CH2)4COOH -> H2N(CH2)6NHCO(CH2)4COOH + H2O} .
\]
The product still carries a free \omterm{def:g11:functional-groups:amine}{amine} group at one end and a free acid
group at the other: each can react with another \omterm{def:g12:polymers:polymer}{monomer}.
\end{solution}

\begin{solution}{exo:g12:polymers:13}
Two \omterm{def:g7:atoms-and-molecules:molecule}{molecules} of water per \omterm{def:g12:polymers:repeat-unit}{repeat unit}. $n = 1000/226.0 =
\qty{4.42}{mol}$ of \omterm{def:g12:polymers:repeat-unit}{repeat units}, so \qty{8.85}{mol} of water:
$8.85 \times 18.0 = \qty{159}{g}$.
\end{solution}

\begin{solution}{exo:g12:polymers:14}
Warm natural rubber is made of chains that slide past one another.
Sulfur bridges link the chains: they stretch but spring back, and no
longer flow. Being cross-linked, a tyre cannot be melted and reshaped.
\end{solution}

\begin{solution}{exo:g12:polymers:15}
$M(\ce{C6H10O5}) = \qty{162.0}{g/mol}$; $n = \num{1.6e6}/162.0 = 9900$
units; length $9900 \times 0.5 = \qty{4900}{nm}$, about
\qty{5}{\micro m}.
\end{solution}

\begin{solution}{pb:g12:polymers:1}
\textbf{1.} Ethane-1,2-diol (two \omterm{def:g11:functional-groups:alcohol}{alcohol} groups) and
benzene-1,4-dicarboxylic acid (two \omterm{def:g11:functional-groups:carboxylic-acid}{carboxylic acid} groups).

\textbf{2.} An \omterm{def:g11:functional-groups:carboxylic-acid}{ester} link; water is released.

\textbf{3.} A \omterm{def:g12:polymers:addition-polymer}{condensation polymer}.

\textbf{4.} $M(\ce{C2H6O2}) = \qty{62.0}{g/mol}$; $M(\ce{C8H6O4}) =
\qty{166.0}{g/mol}$.

\textbf{5.} $62.0 + 166.0 - 2 \times 18.0 = \qty{192.0}{g/mol}$; and
$10 \times 12.0 + 8 \times 1.0 + 4 \times 16.0 = 192.0$.

\textbf{6.} $n = 38400/192.0 = 200$.

\textbf{7.} Two \omterm{def:g11:functional-groups:carboxylic-acid}{ester} links per \omterm{def:g12:polymers:repeat-unit}{repeat unit}: about 400.

\textbf{8.} Its chains are mostly linear, and parts of them line up in
ordered regions.

\textbf{9.} It is a \omterm{def:g8:plastics:thermoplastic}{thermoplastic}: its chains are not cross-linked and
slide past each other when hot.

\textbf{10.} $300/0.80 = \qty{375}{g}$.

\textbf{11.} $375/25 = 15$ bottles.

\textbf{12.} Caps, labels, dirt and scraps are removed or lost during
sorting, washing and spinning.

\textbf{13.} Prevent waste (principle 1): a used object becomes \omterm{def:g5:raw-materials-and-recycling:raw-material}{raw
material}.

\textbf{14.} \ce{C10H8O4 + 2H2O -> C2H6O2 + C8H6O4}.

\textbf{15.} $1000/192.0 = \qty{5.21}{mol}$.

\textbf{16.} Diol: $5.21 \times 62.0 = \qty{323}{g}$; acid:
$5.21 \times 166.0 = \qty{865}{g}$.

\textbf{17.} $2 \times 5.21 \times 18.0 = \qty{188}{g}$ of water;
$1000 + 188 = 323 + 865 = \qty{1188}{g}$: mass is conserved.

\textbf{18.} The \omterm{def:g12:polymers:polymer}{monomers} can be purified and give new PET as good as
new, while melting shortens the chains a little each time.

\textbf{19.} One fleece jacket takes 15 bottles of \qty{25}{g}.
\end{solution}
